\subsection{CT9 --- Overload exhaustion}\label{ct:9}\label{the-charging-principle-in-general-form}
\ctsignature{\(P_{\ast\to9}\to \ctcert{C1};\ \ctint{P_{9\to4},P_{9\to7},P_{9\to8}};\ \ctrec{P_{9\to10}}\).}

CT9 is an exact machine for proving bounded multiplicity by exhausting an
overloaded payer fibre.  Its entry interface fixes the payer, fibre, finite
labels, capacity datum, scope predicate, and C1 claim.  The executable
framework adds the fibre-definition certificate, bounded and overloaded
predicates, the homogeneous-extraction theorem, and exact CT4, CT7, CT8, and
CT10 entry contracts.

The scope node distinguishes three cases.  Absence of any finite labelling is
a typed scope terminal.  A finite but inadequate labelling produces a proof of
coarseness and an exact CT10 refinement input.  A usable labelling produces a
\texttt{ScopedState}.  The \texttt{fibre} node then has one successor and
certifies that the stored fibre is the declared inverse image with its
capacity data.

The overload decision consumes the certified fibre.  Its bounded constructor
returns an exact CT4 input together with the bounded-multiplicity proof; its
overloaded constructor returns an \texttt{OverloadedState}.  The
\texttt{extraction} node has one successor and certifies the homogeneous
subfamily required by the instantiation theorem.  Only then does the routing
node choose C1, an exact CT7 exchange input, or an exact CT8 sequence input.
Consequently CT9 does not assume that an arbitrary labelled family has one of
these outputs: the extraction theorem is an explicit node contract.

\begin{itemize}
\tightlist
\item \textbf{S-Def} is fixed by the payer fibre, finite labels, and capacity datum and certified by \texttt{FibreState}.
\item \textbf{S-Dich} is the bounded/overloaded decision and the exhaustive extraction result.
\item \textbf{S-Comp} is the homogeneous-extraction certificate.
\item \textbf{S-Rout} and \textbf{S-Trig} are carried by exact CT4, CT7, CT8, and CT10 validated inputs.
\end{itemize}

An unlabelled fibre and an undeclared tie-break cannot be represented by a
successful plan.  The scope terminal records failure of finite expressibility;
the CT10 terminal records a concrete finite refinement problem.

\begin{figure}[p]
\centering
\vspace*{-1.75cm}
\resizebox{\textwidth}{!}{%
\begin{tikzpicture}[x=1cm,y=1cm,every node/.style={font=\scriptsize}]
\node[bcwide] (entry) at (0,0) {{\tiny\ttfamily CT9.entry}\\\textbf{Validated payer fibre and capacity}};
\node[bcdec,bclemma,text width=3.55cm,minimum width=4.30cm] (scope) at (0,-2.30) {{\tiny\ttfamily CT9.decide.scope}\\\textbf{Scope, refine, or usable labels?}};
\node[bcscopefail,text width=3.05cm,minimum width=3.05cm] (scopeout) at (7.30,-2.30) {{\tiny\ttfamily CT9.terminal.scope}\\\textbf{Scope exit}};
\node[bcroute,bcrouteneutral,bcdesign,text width=3.05cm,minimum width=3.05cm] (ct10) at (-7.30,-2.30) {{\tiny\ttfamily CT9.terminal.ct10}\\\textbf{CT 10 exit}};
\node[bcwide,bclemma] (fibre) at (0,-4.75) {{\tiny\ttfamily CT9.certify.fibre}\\\textbf{Payer-fibre definition certificate}};
\node[bcdec,bclemma,text width=3.45cm,minimum width=4.20cm] (overload) at (0,-7.15) {{\tiny\ttfamily CT9.decide.overload}\\\textbf{Bounded or overloaded?}};
\node[bcroute,bcrouteneutral,bcoutputroute,text width=3.05cm,minimum width=3.05cm] (ct4) at (7.30,-7.15) {{\tiny\ttfamily CT9.terminal.ct4}\\\textbf{CT 4 exit}};
\node[bcwide,bclemma] (extraction) at (0,-9.65) {{\tiny\ttfamily CT9.certify.extraction}\\\textbf{Homogeneous extraction certificate}};
\node[bcroutechoice,bctheorem,text width=3.25cm,minimum width=4.10cm] (routing) at (0,-12.15) {{\tiny\ttfamily CT9.decide.routing}\\\textbf{Classify extracted object}};
\node[bcclosed,text width=2.85cm,minimum width=2.85cm] (c1) at (-5.20,-15.20) {{\tiny\ttfamily CT9.terminal.c1}\\\textbf{C1}};
\node[bcroute,bcrouteneutral,bcoutputroute,text width=3.05cm,minimum width=3.05cm] (ct7) at (0,-15.20) {{\tiny\ttfamily CT9.terminal.ct7}\\\textbf{CT 7 exit}};
\node[bcroute,bcrouteneutral,bcoutputroute,text width=3.05cm,minimum width=3.05cm] (ct8) at (5.20,-15.20) {{\tiny\ttfamily CT9.terminal.ct8}\\\textbf{CT 8 exit}};
\draw[bcarr] (entry)--(scope);
\draw[bcscopearr] (scope)--node[bclabel,above]{no finite labels}(scopeout);
\draw[bcrecoverarr] (scope)--node[bclabel,above]{labels too coarse}(ct10);
\draw[bcarr] (scope)--node[bclabel,right]{ready}(fibre);
\draw[bcarr] (fibre)--node[bclabel,right]{certified}(overload);
\draw[bchandoff] (overload)--node[bclabel,above]{bounded}(ct4);
\draw[bcarr] (overload)--node[bclabel,right]{overloaded}(extraction);
\draw[bcarr] (extraction)--node[bclabel,right]{certified}(routing);
\draw[bcarr] (routing)--node[bclabel,above,sloped]{target hit}(c1);
\draw[bchandoff] (routing)--node[bclabel,right]{\(P_{9\to7}\)}(ct7);
\draw[bchandoff] (routing)--node[bclabel,above,sloped]{\(P_{9\to8}\)}(ct8);
\end{tikzpicture}%
}
\caption[Closure-tactic 9: exact overload machine]{Closure-tactic 9 as an exact sequential machine.  Fibre definition and homogeneous extraction are one-successor certification boundaries.  Bounded multiplicity returns an exact CT4 input; overload routing occurs only after the extraction theorem has produced its indexed state.}
\label{fig:ct9-overload-exhaustion}
\end{figure}
\FloatBarrier
